\section{Router Intermedio}\label{router-intermedio}

Un'altra possibile alternativa consiste nell'introdurre un router
dedicato (con un'immagine alpine per esempio) che fa da tramite tra il
bridge di verifica e l'interfaccia \passthrough{\lstinline!incusbr0!}.

In questa architettura, i container duplicati vengono collegati a un
bridge L2 isolato. Il router intermedio si aggancia al bridge usando
come IP il gateway originale (\passthrough{\lstinline!10.k.1.1/24!}).

Verso l'esterno, il router ha una seconda interfaccia collegata verso
\passthrough{\lstinline!incusbr0!} da cui ottiene un indirizzo IP.
Vengono applicate le regole di SNAT per mascherare il traffico uscente.
In questo modo la topologia interna della rete di verifica viene
completamente nascosta, evitando i conflitti di instradamento con la
rete di produzione.

Sul piano operativo l'adozione del router richiede di istanziare un
container in più rispetto alla soluzione adottata con OVN.

In OVN inoltre la creazione dello switch, del router logico e delle
regole di NAT è gestita in modo nativo da Incus, mentre l'uso di un
router intermedio impone a BlueAgent di gestirlo manualmente per ogni
sessione di verifica. Questo comporta un lieve aumento della complessità
nel codice e del footprint di memoria.
